% Options for packages loaded elsewhere
% Options for packages loaded elsewhere
\PassOptionsToPackage{unicode}{hyperref}
\PassOptionsToPackage{hyphens}{url}
\PassOptionsToPackage{dvipsnames,svgnames,x11names}{xcolor}
%
\documentclass[
  xelatex]{bxjsarticle}
\usepackage{xcolor}
\usepackage{amsmath,amssymb}
\setcounter{secnumdepth}{-\maxdimen} % remove section numbering
\usepackage{iftex}
\ifPDFTeX
  \usepackage[T1]{fontenc}
  \usepackage[utf8]{inputenc}
  \usepackage{textcomp} % provide euro and other symbols
\else % if luatex or xetex
  \usepackage{unicode-math} % this also loads fontspec
  \defaultfontfeatures{Scale=MatchLowercase}
  \defaultfontfeatures[\rmfamily]{Ligatures=TeX,Scale=1}
\fi
\usepackage{lmodern}
\ifPDFTeX\else
  % xetex/luatex font selection
  \setmainfont[]{Hiragino Mincho ProN}
  \setsansfont[]{Hiragino Sans}
  \setmonofont[]{Hiragino Kaku Gothic ProN}
\fi
% Use upquote if available, for straight quotes in verbatim environments
\IfFileExists{upquote.sty}{\usepackage{upquote}}{}
\IfFileExists{microtype.sty}{% use microtype if available
  \usepackage[]{microtype}
  \UseMicrotypeSet[protrusion]{basicmath} % disable protrusion for tt fonts
}{}
\makeatletter
\@ifundefined{KOMAClassName}{% if non-KOMA class
  \IfFileExists{parskip.sty}{%
    \usepackage{parskip}
  }{% else
    \setlength{\parindent}{0pt}
    \setlength{\parskip}{6pt plus 2pt minus 1pt}}
}{% if KOMA class
  \KOMAoptions{parskip=half}}
\makeatother
% Make \paragraph and \subparagraph free-standing
\makeatletter
\ifx\paragraph\undefined\else
  \let\oldparagraph\paragraph
  \renewcommand{\paragraph}{
    \@ifstar
      \xxxParagraphStar
      \xxxParagraphNoStar
  }
  \newcommand{\xxxParagraphStar}[1]{\oldparagraph*{#1}\mbox{}}
  \newcommand{\xxxParagraphNoStar}[1]{\oldparagraph{#1}\mbox{}}
\fi
\ifx\subparagraph\undefined\else
  \let\oldsubparagraph\subparagraph
  \renewcommand{\subparagraph}{
    \@ifstar
      \xxxSubParagraphStar
      \xxxSubParagraphNoStar
  }
  \newcommand{\xxxSubParagraphStar}[1]{\oldsubparagraph*{#1}\mbox{}}
  \newcommand{\xxxSubParagraphNoStar}[1]{\oldsubparagraph{#1}\mbox{}}
\fi
\makeatother

\usepackage{color}
\usepackage{fancyvrb}
\newcommand{\VerbBar}{|}
\newcommand{\VERB}{\Verb[commandchars=\\\{\}]}
\DefineVerbatimEnvironment{Highlighting}{Verbatim}{commandchars=\\\{\}}
% Add ',fontsize=\small' for more characters per line
\usepackage{framed}
\definecolor{shadecolor}{RGB}{241,243,245}
\newenvironment{Shaded}{\begin{snugshade}}{\end{snugshade}}
\newcommand{\AlertTok}[1]{\textcolor[rgb]{0.68,0.00,0.00}{#1}}
\newcommand{\AnnotationTok}[1]{\textcolor[rgb]{0.37,0.37,0.37}{#1}}
\newcommand{\AttributeTok}[1]{\textcolor[rgb]{0.40,0.45,0.13}{#1}}
\newcommand{\BaseNTok}[1]{\textcolor[rgb]{0.68,0.00,0.00}{#1}}
\newcommand{\BuiltInTok}[1]{\textcolor[rgb]{0.00,0.23,0.31}{#1}}
\newcommand{\CharTok}[1]{\textcolor[rgb]{0.13,0.47,0.30}{#1}}
\newcommand{\CommentTok}[1]{\textcolor[rgb]{0.37,0.37,0.37}{#1}}
\newcommand{\CommentVarTok}[1]{\textcolor[rgb]{0.37,0.37,0.37}{\textit{#1}}}
\newcommand{\ConstantTok}[1]{\textcolor[rgb]{0.56,0.35,0.01}{#1}}
\newcommand{\ControlFlowTok}[1]{\textcolor[rgb]{0.00,0.23,0.31}{\textbf{#1}}}
\newcommand{\DataTypeTok}[1]{\textcolor[rgb]{0.68,0.00,0.00}{#1}}
\newcommand{\DecValTok}[1]{\textcolor[rgb]{0.68,0.00,0.00}{#1}}
\newcommand{\DocumentationTok}[1]{\textcolor[rgb]{0.37,0.37,0.37}{\textit{#1}}}
\newcommand{\ErrorTok}[1]{\textcolor[rgb]{0.68,0.00,0.00}{#1}}
\newcommand{\ExtensionTok}[1]{\textcolor[rgb]{0.00,0.23,0.31}{#1}}
\newcommand{\FloatTok}[1]{\textcolor[rgb]{0.68,0.00,0.00}{#1}}
\newcommand{\FunctionTok}[1]{\textcolor[rgb]{0.28,0.35,0.67}{#1}}
\newcommand{\ImportTok}[1]{\textcolor[rgb]{0.00,0.46,0.62}{#1}}
\newcommand{\InformationTok}[1]{\textcolor[rgb]{0.37,0.37,0.37}{#1}}
\newcommand{\KeywordTok}[1]{\textcolor[rgb]{0.00,0.23,0.31}{\textbf{#1}}}
\newcommand{\NormalTok}[1]{\textcolor[rgb]{0.00,0.23,0.31}{#1}}
\newcommand{\OperatorTok}[1]{\textcolor[rgb]{0.37,0.37,0.37}{#1}}
\newcommand{\OtherTok}[1]{\textcolor[rgb]{0.00,0.23,0.31}{#1}}
\newcommand{\PreprocessorTok}[1]{\textcolor[rgb]{0.68,0.00,0.00}{#1}}
\newcommand{\RegionMarkerTok}[1]{\textcolor[rgb]{0.00,0.23,0.31}{#1}}
\newcommand{\SpecialCharTok}[1]{\textcolor[rgb]{0.37,0.37,0.37}{#1}}
\newcommand{\SpecialStringTok}[1]{\textcolor[rgb]{0.13,0.47,0.30}{#1}}
\newcommand{\StringTok}[1]{\textcolor[rgb]{0.13,0.47,0.30}{#1}}
\newcommand{\VariableTok}[1]{\textcolor[rgb]{0.07,0.07,0.07}{#1}}
\newcommand{\VerbatimStringTok}[1]{\textcolor[rgb]{0.13,0.47,0.30}{#1}}
\newcommand{\WarningTok}[1]{\textcolor[rgb]{0.37,0.37,0.37}{\textit{#1}}}

\usepackage{longtable,booktabs,array}
\usepackage{calc} % for calculating minipage widths
% Correct order of tables after \paragraph or \subparagraph
\usepackage{etoolbox}
\makeatletter
\patchcmd\longtable{\par}{\if@noskipsec\mbox{}\fi\par}{}{}
\makeatother
% Allow footnotes in longtable head/foot
\IfFileExists{footnotehyper.sty}{\usepackage{footnotehyper}}{\usepackage{footnote}}
\makesavenoteenv{longtable}
\usepackage{graphicx}
\makeatletter
\newsavebox\pandoc@box
\newcommand*\pandocbounded[1]{% scales image to fit in text height/width
  \sbox\pandoc@box{#1}%
  \Gscale@div\@tempa{\textheight}{\dimexpr\ht\pandoc@box+\dp\pandoc@box\relax}%
  \Gscale@div\@tempb{\linewidth}{\wd\pandoc@box}%
  \ifdim\@tempb\p@<\@tempa\p@\let\@tempa\@tempb\fi% select the smaller of both
  \ifdim\@tempa\p@<\p@\scalebox{\@tempa}{\usebox\pandoc@box}%
  \else\usebox{\pandoc@box}%
  \fi%
}
% Set default figure placement to htbp
\def\fps@figure{htbp}
\makeatother





\setlength{\emergencystretch}{3em} % prevent overfull lines

\providecommand{\tightlist}{%
  \setlength{\itemsep}{0pt}\setlength{\parskip}{0pt}}



 


\makeatletter
\@ifpackageloaded{caption}{}{\usepackage{caption}}
\AtBeginDocument{%
\ifdefined\contentsname
  \renewcommand*\contentsname{Table of contents}
\else
  \newcommand\contentsname{Table of contents}
\fi
\ifdefined\listfigurename
  \renewcommand*\listfigurename{List of Figures}
\else
  \newcommand\listfigurename{List of Figures}
\fi
\ifdefined\listtablename
  \renewcommand*\listtablename{List of Tables}
\else
  \newcommand\listtablename{List of Tables}
\fi
\ifdefined\figurename
  \renewcommand*\figurename{Figure}
\else
  \newcommand\figurename{Figure}
\fi
\ifdefined\tablename
  \renewcommand*\tablename{Table}
\else
  \newcommand\tablename{Table}
\fi
}
\@ifpackageloaded{float}{}{\usepackage{float}}
\floatstyle{ruled}
\@ifundefined{c@chapter}{\newfloat{codelisting}{h}{lop}}{\newfloat{codelisting}{h}{lop}[chapter]}
\floatname{codelisting}{Listing}
\newcommand*\listoflistings{\listof{codelisting}{List of Listings}}
\makeatother
\makeatletter
\makeatother
\makeatletter
\@ifpackageloaded{caption}{}{\usepackage{caption}}
\@ifpackageloaded{subcaption}{}{\usepackage{subcaption}}
\makeatother
\usepackage{bookmark}
\IfFileExists{xurl.sty}{\usepackage{xurl}}{} % add URL line breaks if available
\urlstyle{same}
\hypersetup{
  pdftitle={春学期ドラフト},
  colorlinks=true,
  linkcolor={blue},
  filecolor={Maroon},
  citecolor={Blue},
  urlcolor={Blue},
  pdfcreator={LaTeX via pandoc}}


\title{春学期ドラフト}
\author{}
\date{}
\begin{document}
\maketitle


\subsection{概要}\label{ux6982ux8981}

本研究の目的：秋学期に行う世論調査に向け、その世論調査の実験意義を明確にする。先行研究(trust
in
government)では、政府の信頼が科学技術に対して、大きな影響を与えると述べられていたが、それは実験実施国であったエストニアの影響ではないか。オンライン選挙が導入されているという前提があるから、実験用の仮想の国に対してもそのような結果が出たのではないか。では、それがオンライン選挙が導入されていない日本という国において、同じような結果が出るのか。それを確かめるために、秋学期に世論調査実験を行う。今学期は、その世論調査の実験意義を明確にするために、既存の世論調査を用いて、政府への信頼が科学技術への信頼に与える影響を検証する。
秋学期に行う実験内容：オンライン投票(i-voting)において政府への信頼がどのような影響を与えるのか、先行研究のレプリケーションのようなものを題材とする。（これは、本分析の世論調査項目の内容の中でも、情報通信技術や科学技術イノベーションに着目する理由となることを明記して欲しい。）
本研究は、政府が推進する特定分野の政策に対する国民の重要度認識において、「政府への信頼」がどのような役割を果たすかを検証することを目的とする。内閣府実施の「くらしと科学技術に関する意識調査」の個票データを用い、ロジスティック回帰分析を行った。分析の結果、個人の属性（性別・年齢）の影響を統制した後でも、政府への信頼度は「情報通信政策」および「科学技術イノベーション政策」の重要度認識に対して、一貫して統計的に極めて有意な正の影響を与えることが明らかになった。また、統制変数の影響は政策分野によって異なり、科学技術政策では性差が、情報通信政策では特定の年齢層の影響が見られるなど、示唆に富む知見が得られた。これらの結果は、政策への支持形成における信頼の普遍的な重要性と、分野ごとに異なる考慮事項が存在することを示している。

\subsection{1.
イントロダクション}\label{ux30a4ux30f3ux30c8ux30edux30c0ux30afux30b7ux30e7ux30f3}

情報通信技術や科学技術イノベーションは、現代国家の競争力や国民生活の質を左右する重要な政策分野である。これらの政策を効果的に推進するには、国民の理解と支持が不可欠となる。しかし、政策の内容が高度に専門化・複雑化する中で、国民一人ひとりが政策の詳細を理解し、その重要性を判断することは容易ではない。

このような状況下で、人々は政策判断の際に何らかの簡便な判断基準（ヒューリスティクス）を用いると考えられる。その中でも、政策の主体である「政府」という組織への信頼は、最も重要な判断基準の一つとなり得るだろう。

本研究では、この「政府への信頼」という要素に着目し、「政府への信頼度は、情報通信や科学技術といった特定分野の政策に対する国民の重要度認識にどのような影響を与えるのか？」というリサーチ・クエスチョンを立て、実証的なデータ分析を通じてその関係性を明らかにすることを試みる。

\subsection{2. 文献レビュー}\label{ux6587ux732eux30ecux30d3ux30e5ux30fc}

政策受容や技術受容に関する研究において、市民の「信頼（Trust）」は中心的な役割を果たす変数として広く研究されてきた。特に、電子投票（i-voting）のような新しい技術の導入に関する先行研究（e.g.,
Ebbers, Pieterson, \& Noordman,
2008）では、技術そのものへの信頼と同じく、あるいはそれ以上に、運営主体である政府への信頼が受容意図を左右することが示されている。

これらの研究は、人々が対象に関する十分な知識を持たない場合、その対象を管理・推進する組織への信頼を代理的な判断基準として用いることを示唆している。本研究は、この議論を電子投票という特定の文脈から、より一般的な「情報通信政策」および「科学技術イノベーション政策」全体へと拡張し、日本社会の文脈でその妥当性を検証するものである。

\subsection{3. 理論・仮説}\label{ux7406ux8ad6ux4eeeux8aac}

\begin{itemize}
\tightlist
\item
  \textbf{理論}
\end{itemize}

本研究の理論的枠組みは、ヒューリスティック理論に基づく。複雑な政策課題に直面した際、人々は全ての情報を合理的に処理するのではなく、認知的な負荷を軽減するためのショートカットを用いる。その代表的なものが「信頼ヒューリスティック」である。
政策の推進主体である政府を信頼している個人は、政府が提案する政策もまた、国民にとって有益で重要なものであると推論する傾向があると考えられる。逆に、政府に不信感を抱いている個人は、政府の政策に対しても懐疑的、あるいは無関心になりやすいだろう。
この理論に基づき、本研究では以下の仮説を立てる。

\begin{itemize}
\tightlist
\item
  \textbf{仮説}
\end{itemize}

政府への信頼度が高い人ほど、政府が推進する政策（情報通信政策、科学技術イノベーション政策）を重要だと認識する傾向がある。

\subsection{4.
データ/関連概念の測定方法}\label{ux30c7ux30fcux30bfux95a2ux9023ux6982ux5ff5ux306eux6e2cux5b9aux65b9ux6cd5}

\subsubsection{データ}\label{ux30c7ux30fcux30bf}

本研究では、内閣府が実施した「くらしと科学技術に関する意識調査」の個票データ（\texttt{RawData\_OpinionPoll\_STIpolicy.csv}）を使用する。

\subsubsection{変数の操作的定義}\label{ux5909ux6570ux306eux64cdux4f5cux7684ux5b9aux7fa9}

分析に使用するため、調査票の回答を以下のように処理し、変数を定義した。

\begin{itemize}
\tightlist
\item
  \textbf{従属変数}:

  \begin{itemize}
  \tightlist
  \item
    \textbf{情報通信政策の重要度認識 (\texttt{is\_important\_q8\_5})}:
    設問Q8.5「情報通信政策」について「重要」または「まあ重要」と回答した場合を\texttt{1}、それ以外を\texttt{0}とした。
  \item
    \textbf{科学技術政策の重要度認識 (\texttt{is\_important\_q8\_19})}:
    設問Q8.19「科学技術イノベーション政策」について「重要」または「まあ重要」と回答した場合を\texttt{1}、それ以外を\texttt{0}とした。
  \end{itemize}
\item
  \textbf{独立変数}:

  \begin{itemize}
  \tightlist
  \item
    \textbf{政府への信頼 (\texttt{trust\_government})}:
    設問Q11.8（先行研究との接続性を考慮し選択）について「信頼できる」または「まあ信頼できる」と回答した場合を\texttt{1}、それ以外を\texttt{0}とした。
  \end{itemize}
\item
  \textbf{統制変数}:

  \begin{itemize}
  \tightlist
  \item
    \textbf{性別 (\texttt{gender})}:
    設問F1の回答を「男性」「女性」のラベルを持つ因子（factor）として定義した。
  \item
    \textbf{年齢 (\texttt{age\_category})}:
    設問F2の14の年齢カテゴリ（「16〜19歳」など）を、それぞれラベルを持つ因子（factor）として定義した。
  \end{itemize}
\end{itemize}

\subsection{5. 分析}\label{ux5206ux6790}

上記の変数を用いて、従属変数が二値（0か1）であることからロジスティック回帰分析を行った。

\subsubsection{モデル1:
単純モデル（統制なし）}\label{ux30e2ux30c7ux30eb1-ux5358ux7d14ux30e2ux30c7ux30ebux7d71ux5236ux306aux3057}

\begin{Shaded}
\begin{Highlighting}[]
\ControlFlowTok{if}\NormalTok{ (}\SpecialCharTok{!}\FunctionTok{requireNamespace}\NormalTok{(}\StringTok{"tidyverse"}\NormalTok{, }\AttributeTok{quietly =} \ConstantTok{TRUE}\NormalTok{)) \{}
  \FunctionTok{install.packages}\NormalTok{(}\StringTok{"tidyverse"}\NormalTok{, }\AttributeTok{repos =} \StringTok{"https://cran.rstudio.com/"}\NormalTok{)}
\NormalTok{\}}
\ControlFlowTok{if}\NormalTok{ (}\SpecialCharTok{!}\FunctionTok{requireNamespace}\NormalTok{(}\StringTok{"broom"}\NormalTok{, }\AttributeTok{quietly =} \ConstantTok{TRUE}\NormalTok{)) \{}
  \FunctionTok{install.packages}\NormalTok{(}\StringTok{"broom"}\NormalTok{, }\AttributeTok{repos =} \StringTok{"https://cran.rstudio.com/"}\NormalTok{)}
\NormalTok{\}}
\FunctionTok{library}\NormalTok{(tidyverse)}
\end{Highlighting}
\end{Shaded}

\begin{verbatim}
Warning: パッケージ 'lubridate' はバージョン 4.4.1 の R の下で造られました
\end{verbatim}

\begin{verbatim}
-- Attaching core tidyverse packages ------------------------ tidyverse 2.0.0 --
v dplyr     1.1.4     v readr     2.1.5
v forcats   1.0.0     v stringr   1.5.1
v ggplot2   3.5.1     v tibble    3.2.1
v lubridate 1.9.4     v tidyr     1.3.1
v purrr     1.0.2     
-- Conflicts ------------------------------------------ tidyverse_conflicts() --
x dplyr::filter() masks stats::filter()
x dplyr::lag()    masks stats::lag()
i Use the conflicted package (<http://conflicted.r-lib.org/>) to force all conflicts to become errors
\end{verbatim}

\begin{Shaded}
\begin{Highlighting}[]
\FunctionTok{library}\NormalTok{(broom)}
\end{Highlighting}
\end{Shaded}

\begin{verbatim}
Warning: パッケージ 'broom' はバージョン 4.4.1 の R の下で造られました
\end{verbatim}

\begin{Shaded}
\begin{Highlighting}[]
\NormalTok{current\_dir }\OtherTok{\textless{}{-}} \FunctionTok{setwd}\NormalTok{(}\StringTok{"/Users/karubeshougo/Uni/seminar/graduation\_thesis"}\NormalTok{)}
\NormalTok{origin\_data }\OtherTok{\textless{}{-}} \FunctionTok{read.csv}\NormalTok{(}\StringTok{"Statics/くらしと科学技術に関する意識調査/sources/origin\_data/RawData\_OpinionPoll\_STIpolicy.csv"}\NormalTok{, }\AttributeTok{header =} \ConstantTok{TRUE}\NormalTok{, }\AttributeTok{sep =} \StringTok{","}\NormalTok{, }\AttributeTok{quote =} \StringTok{"}\SpecialCharTok{\textbackslash{}"}\StringTok{"}\NormalTok{, }\AttributeTok{dec =} \StringTok{"."}\NormalTok{, }\AttributeTok{fill =} \ConstantTok{TRUE}\NormalTok{, }\AttributeTok{comment.char =} \StringTok{""}\NormalTok{, }\AttributeTok{fileEncoding =} \StringTok{"CP932"}\NormalTok{, }\AttributeTok{encoding =} \StringTok{"UTF{-}8"}\NormalTok{)}

\NormalTok{analysis\_data }\OtherTok{\textless{}{-}} \FunctionTok{data.frame}\NormalTok{(}
\CommentTok{\# [従属変数]q8.5, q8.19 において、重要であると考えている}
    \AttributeTok{is\_important\_q8\_5 =} \FunctionTok{ifelse}\NormalTok{(origin\_data}\SpecialCharTok{$}\NormalTok{Q8\_5 }\SpecialCharTok{==} \DecValTok{1} \SpecialCharTok{|}\NormalTok{ origin\_data}\SpecialCharTok{$}\NormalTok{Q8\_5 }\SpecialCharTok{==} \DecValTok{2}\NormalTok{, }\DecValTok{1}\NormalTok{, }\DecValTok{0}\NormalTok{),}
    \AttributeTok{is\_important\_q8\_19 =} \FunctionTok{ifelse}\NormalTok{(origin\_data}\SpecialCharTok{$}\NormalTok{Q8\_19 }\SpecialCharTok{==} \DecValTok{1} \SpecialCharTok{|}\NormalTok{ origin\_data}\SpecialCharTok{$}\NormalTok{Q8\_19 }\SpecialCharTok{==} \DecValTok{2}\NormalTok{, }\DecValTok{1}\NormalTok{, }\DecValTok{0}\NormalTok{),}
    
\CommentTok{\# [独立変数] 政府への信頼度}
    \AttributeTok{trust\_government =} \FunctionTok{ifelse}\NormalTok{(origin\_data}\SpecialCharTok{$}\NormalTok{Q11\_8 }\SpecialCharTok{==} \DecValTok{1} \SpecialCharTok{|}\NormalTok{ origin\_data}\SpecialCharTok{$}\NormalTok{Q11\_8 }\SpecialCharTok{==} \DecValTok{2}\NormalTok{, }\DecValTok{1}\NormalTok{, }\DecValTok{0}\NormalTok{)}
\NormalTok{)}

\CommentTok{\# モデル1}
\CommentTok{\# 政府への信頼が「情報通信政策」の重要度認識に与える影響}
\FunctionTok{print}\NormalTok{(}\StringTok{"情報通信政策の重要度認識における政府への信頼の影響"}\NormalTok{)}
\end{Highlighting}
\end{Shaded}

\begin{verbatim}
[1] "情報通信政策の重要度認識における政府への信頼の影響"
\end{verbatim}

\begin{Shaded}
\begin{Highlighting}[]
\NormalTok{model\_1 }\OtherTok{\textless{}{-}} \FunctionTok{glm}\NormalTok{(is\_important\_q8\_5 }\SpecialCharTok{\textasciitilde{}}\NormalTok{ trust\_government, }\AttributeTok{data =}\NormalTok{ analysis\_data, }\AttributeTok{family =} \StringTok{"binomial"}\NormalTok{)}
\FunctionTok{summary}\NormalTok{(model\_1)}
\end{Highlighting}
\end{Shaded}

\begin{verbatim}

Call:
glm(formula = is_important_q8_5 ~ trust_government, family = "binomial", 
    data = analysis_data)

Coefficients:
                 Estimate Std. Error z value Pr(>|z|)    
(Intercept)       0.66989    0.09321   7.187 6.63e-13 ***
trust_government  0.56572    0.15509   3.648 0.000265 ***
---
Signif. codes:  0 '***' 0.001 '**' 0.01 '*' 0.05 '.' 0.1 ' ' 1

(Dispersion parameter for binomial family taken to be 1)

    Null deviance: 1069.6  on 886  degrees of freedom
Residual deviance: 1055.9  on 885  degrees of freedom
AIC: 1059.9

Number of Fisher Scoring iterations: 4
\end{verbatim}

\begin{Shaded}
\begin{Highlighting}[]
\CommentTok{\# モデル2}
\CommentTok{\# 政府への信頼が「科学技術イノベーション政策」の重要度認識に与える影響}
\FunctionTok{print}\NormalTok{(}\StringTok{"科学技術イノベーション政策の重要度認識における政府への信頼の影響"}\NormalTok{)}
\end{Highlighting}
\end{Shaded}

\begin{verbatim}
[1] "科学技術イノベーション政策の重要度認識における政府への信頼の影響"
\end{verbatim}

\begin{Shaded}
\begin{Highlighting}[]
\NormalTok{model\_2 }\OtherTok{\textless{}{-}} \FunctionTok{glm}\NormalTok{(is\_important\_q8\_19 }\SpecialCharTok{\textasciitilde{}}\NormalTok{ trust\_government, }\AttributeTok{data =}\NormalTok{ analysis\_data, }\AttributeTok{family =} \StringTok{"binomial"}\NormalTok{)}
\FunctionTok{summary}\NormalTok{(model\_2)}
\end{Highlighting}
\end{Shaded}

\begin{verbatim}

Call:
glm(formula = is_important_q8_19 ~ trust_government, family = "binomial", 
    data = analysis_data)

Coefficients:
                 Estimate Std. Error z value Pr(>|z|)    
(Intercept)       0.50046    0.09099   5.500 3.80e-08 ***
trust_government  0.78184    0.15508   5.042 4.62e-07 ***
---
Signif. codes:  0 '***' 0.001 '**' 0.01 '*' 0.05 '.' 0.1 ' ' 1

(Dispersion parameter for binomial family taken to be 1)

    Null deviance: 1098.3  on 886  degrees of freedom
Residual deviance: 1071.7  on 885  degrees of freedom
AIC: 1075.7

Number of Fisher Scoring iterations: 4
\end{verbatim}

まず、政府への信頼が各政策の重要度認識に与える直接的な影響を見るため、統制変数を投入しない単純なモデルを構築した。結果、\textbf{両政策において、政府への信頼は統計的に極めて有意な（p
\textless{}
0.001）正の影響}を与えていることが確認された。これは仮説を支持する基本的な結果である。

\subsubsection{モデル2:
多変量モデル（統制あり）}\label{ux30e2ux30c7ux30eb2-ux591aux5909ux91cfux30e2ux30c7ux30ebux7d71ux5236ux3042ux308a}

次に、性別と年齢の影響を取り除くため、これらを統制変数として投入した重回帰ロジスティック分析を行った。

\paragraph{主要な発見}\label{ux4e3bux8981ux306aux767aux898b}

最大の発見は、\textbf{性別と年齢の影響を考慮してもなお、\texttt{trust\_government}（政府への信頼）の係数は両モデルにおいて統計的に極めて有意なままであり、その効果が非常に頑健（ロバスト）である}ことが証明された点である。

\paragraph{統制変数の効果}\label{ux7d71ux5236ux5909ux6570ux306eux52b9ux679c}

統制変数の効果は、政策分野によって異なる興味深いパターンを示した。

\begin{itemize}
\tightlist
\item
  \textbf{性別の影響}:

  \begin{itemize}
  \tightlist
  \item
    \textbf{情報通信政策}: 性別による有意な差は見られなかった。
  \item
    \textbf{科学技術政策}: 統計的に有意な差が見られ（p \textless{}
    0.05）、男性に比べて女性の方が重要だと考える傾向がやや弱いことが示された。
  \end{itemize}
\end{itemize}

\begin{itemize}
\tightlist
\item
  \textbf{年齢の影響}:

  \begin{itemize}
  \tightlist
  \item
    \textbf{情報通信政策}:
    16〜19歳を基準とした場合、\textbf{50代（50〜59歳）が他の世代に比べて統計的に有意に重要だと考える}傾向が強い、山なりのパターンが見られた。
  \item
    \textbf{科学技術政策}:
    多くの年齢層で基準との間に有意な差は見られず、特定の年齢層の影響は限定的であった。
  \end{itemize}
\end{itemize}

これらの結果は、政府への信頼という普遍的な影響に加え、政策テーマごとに異なる人口統計学的属性が関連することを示唆している。

\subsection{6. 結論}\label{ux7d50ux8ad6}

本分析により、以下の3点が明らかになった。

\begin{enumerate}
\def\labelenumi{\arabic{enumi}.}
\tightlist
\item
  \textbf{政府への信頼は、政策分野や個人の属性を超えて、政策の重要度認識に影響を与える最も強力な共通因子である。}
\item
  \textbf{性差の影響は、科学技術政策という特定の分野において見られる。}
\item
  \textbf{年齢の影響は、情報通信政策において特に50代の関心の高さとして現れる。}
\end{enumerate}

本研究の\textbf{含意}として、政府がハイテク政策への国民の支持を醸成する上で、小手先の政策広報以上に、政府そのものへの信頼を地道に構築することが根源的に重要であることが示された。また、政策分野ごとにコミュニケーションの対象とすべき層が異なる可能性も示唆された。

本研究の\textbf{限界}として、クロスセクションデータを用いているため、変数間の因果関係を断定することはできない点が挙げられる。また、使用した信頼の尺度は一設問のみであり、より多角的な指標を用いることで、さらに豊かな知見が得られる可能性がある。今後の研究では、パネルデータや実験的手法を用いた、より厳密な因果関係の検証が期待される。

\subsection{参考文献}\label{ux53c2ux8003ux6587ux732e}

Ebbers, W. E., Pieterson, W. J., \& Noordman, H. N. (2008). Trust in
Government or in Technology? What Really Drives Internet Voting. In
\emph{Electronic Government} (pp.~50-61). Springer.
くらしと科学技術に関する意識調査：https://www.nistep.go.jp/wp/wp-content/uploads/Questionnire\_ja.pdf




\end{document}
